% ============================================================
% research_report.tex
% Generalizable Deepfake Speech Detection Using SVMs
% with Calibrated Probabilistic Inference
% Course: 22MAT220 – Mathematics for Computing III
% ============================================================

\documentclass[conference]{IEEEtran}
\IEEEoverridecommandlockouts

\usepackage{amsmath,amssymb,amsfonts}
\usepackage{algorithmic}
\usepackage{graphicx}
\usepackage{textcomp}
\usepackage{xcolor}
\usepackage{booktabs}
\usepackage{url}
\usepackage{hyperref}
\usepackage{array}
\usepackage{cite}

\hypersetup{colorlinks=true, linkcolor=blue, citecolor=blue, urlcolor=blue}

% ─────────────────────────────────────────────
\begin{document}

\title{Generalizable Deepfake Speech Detection Using Support Vector Machines
with Calibrated Probabilistic Inference}

\author{
  \IEEEauthorblockN{[Author 1 Name] \quad [Roll Number]}
  \IEEEauthorblockA{Amrita School of Engineering, Bengaluru\\
  Department of Artificial Intelligence and Data Science\\
  \texttt{[email@blr.amrita.edu]}}
  \and
  \IEEEauthorblockN{[Author 2 Name] \quad [Roll Number]}
  \IEEEauthorblockA{Amrita School of Engineering, Bengaluru\\
  Department of Artificial Intelligence and Data Science\\
  \texttt{[email@blr.amrita.edu]}}
  \and
  \IEEEauthorblockN{Supervisor: [Faculty Name]}
  \IEEEauthorblockA{Department of Mathematics\\
  Amrita School of Engineering, Bengaluru}
}

\maketitle

% ─────────────────────────────────────────────
\begin{abstract}
The rapid proliferation of AI-generated speech synthesis has created an urgent need
for automatic deepfake audio detection systems. This paper presents a complete,
mathematically rigorous pipeline for binary classification of speech recordings as
\emph{real} (authentic human speech) or \emph{fake} (AI-generated/converted speech),
using Support Vector Machines (SVMs) trained on handcrafted acoustic features.
Our implementation includes a Linear SVM derived from first principles using NumPy
and mini-batch subgradient descent with convergence analysis, alongside scikit-learn
baselines (LinearSVC, SVC with RBF kernel, Logistic Regression, Random Forest).
Features include Mel-Frequency Cepstral Coefficients (MFCC), Linear-Frequency Cepstral
Coefficients (LFCC), and spectral/prosodic statistics, with a systematic ablation study.
Probability estimates are obtained via Platt scaling on a held-out calibration set.
The system is evaluated on the ASVspoof~2019 Logical Access benchmark and an
independent external dataset to assess generalization beyond training conditions.
An interactive Gradio-based web application (EchoVerify AI) allows faculty or end-users
to upload arbitrary audio files and receive calibrated REAL/FAKE probability estimates
in real time. We report all evaluation metrics honestly, including the generalization
gap between internal and external performance, and explicitly document limitations.
\end{abstract}

\begin{IEEEkeywords}
deepfake audio detection, support vector machines, hinge loss, Platt scaling,
MFCC, LFCC, ASVspoof, anti-spoofing, subgradient descent, calibrated probabilities
\end{IEEEkeywords}

% ─────────────────────────────────────────────
\section{Introduction}
\label{sec:intro}

Voice synthesis technology has advanced to the point where AI-generated speech is
perceptually indistinguishable from authentic human speech \cite{wavenet2016, tacotron2017}.
Text-to-speech (TTS) systems based on neural vocoders and transformer architectures
can clone a speaker's voice from a few seconds of reference audio \cite{valle2023}.
This creates serious risks: audio deepfakes have been used in fraud, misinformation,
and identity impersonation.

Automatic Speaker Verification (ASV) systems~\cite{asvspoof2019} are increasingly
targeted by such attacks. The ASVspoof challenge series provides a standardized
benchmark for Countermeasure (CM) development. Existing state-of-the-art detectors
use deep learning (Wav2Vec, RawBoost, AASIST~\cite{jung2022aasist}) and achieve near-zero
equal error rates on in-domain data, but suffer significant degradation on novel generators.

This project takes a principled approach rooted in the objectives of the course
22MAT220 (Mathematics for Computing~III): we implement the complete mathematical
machinery of SVMs from scratch, provide convergence analysis, and deploy a fully
functional interactive inference system. While neural approaches are discussed as
future work, the core contribution is an academically defensible SVM-based detector
with calibrated probabilistic output.

\subsection{Contributions}
\begin{enumerate}
  \item A NumPy SVM implementation from first principles with subgradient analysis.
  \item A comparative study across five model families.
  \item A feature ablation study (MFCC, LFCC, spectral, prosodic).
  \item Platt-calibrated probability estimates with ECE analysis.
  \item External generalization evaluation with gap analysis.
  \item An interactive web application (EchoVerify AI) launchable from Google Colab.
\end{enumerate}

% ─────────────────────────────────────────────
\section{Literature Review}
\label{sec:literature}

Sahidullah et al.~\cite{sahidullah2015} established MFCC and LFCC baselines for
ASVspoof~2015, demonstrating that linear-frequency features outperform mel-scale
features for detecting TTS attacks because the linear scale preserves high-frequency
vocoder artefacts. This insight motivates our combined MFCC+LFCC feature set.

The ASVspoof~2019~\cite{asvspoof2019} challenge introduced 19 attack types and showed
that Gaussian Mixture Model (GMM) classifiers using LFCCs achieve competitive baseline
performance, with Equal Error Rates under 5\% on the development set. However, these
baselines degrade significantly on the evaluation set containing new generators,
confirming the generalization challenge.

Todisco et al.~\cite{todisco2016} proposed Constant-Q Cepstral Coefficients (CQCCs)
as an alternative to MFCC and LFCC, achieving better results on some attack types.
Our work focuses on MFCC and LFCC as well-understood baselines.

For probability calibration, Platt~\cite{platt1999} showed that sigmoid fitting
on held-out scores produces well-calibrated probabilities for SVM outputs. The
alternative isotonic regression calibration~\cite{niculescu2005} requires more
calibration data but makes no parametric assumption. We use Platt scaling.

Recent deep learning approaches such as AASIST~\cite{jung2022aasist} use graph
attention networks on raw waveforms, achieving $<1\%$ EER on ASVspoof~2019 LA eval.
However, these methods require large datasets and are less interpretable. For this
academic project, we favour interpretability and mathematical rigor over raw accuracy.

Generalization to in-the-wild deepfakes was studied by M\"{u}ller et al.~\cite{muller2022},
who showed significant performance drops for models trained only on ASVspoof data
when tested on real-world deepfakes. Our external evaluation follows this methodology.

% ─────────────────────────────────────────────
\section{Mathematical Formulation}
\label{sec:math}

\subsection{Decision Hyperplane}

Let $\mathbf{x}_i \in \mathbb{R}^d$ be the feature vector for sample $i$ and
$y_i \in \{-1, +1\}$ be the class label ($+1$ = FAKE, $-1$ = REAL). The SVM
finds a hyperplane:
\begin{equation}
  f(\mathbf{x}) = \mathbf{w}^\top \mathbf{x} + b
\end{equation}
with decision rule $\hat{y} = \text{sign}(f(\mathbf{x}))$.

\subsection{Soft-Margin Primal Formulation}

\begin{equation}
  \min_{\mathbf{w}, b, \boldsymbol{\xi}} \quad
  \frac{1}{2}\|\mathbf{w}\|^2 + \frac{C}{n}\sum_{i=1}^n \xi_i
\end{equation}
\begin{equation}
  \text{subject to: } y_i(\mathbf{w}^\top \mathbf{x}_i + b) \geq 1 - \xi_i,
  \quad \xi_i \geq 0
\end{equation}

Equivalently, using the hinge loss $\ell_H(t) = \max(0, 1-t)$:
\begin{equation}
  \min_{\mathbf{w}, b} \quad
  \frac{1}{2}\|\mathbf{w}\|^2 + \frac{C}{n}\sum_{i=1}^n
  \max\!\left(0,\; 1 - y_i(\mathbf{w}^\top \mathbf{x}_i + b)\right)
  \label{eq:primal}
\end{equation}

\subsection{Subgradient Computation}

Let $m_i = y_i f(\mathbf{x}_i)$. The subgradient of~\eqref{eq:primal} w.r.t. $\mathbf{w}$:

\begin{equation}
  \partial_\mathbf{w} J =
  \begin{cases}
    \mathbf{w} & \text{if } m_i \geq 1 \\
    \mathbf{w} - \dfrac{C}{n}\,y_i\,\mathbf{x}_i & \text{if } m_i < 1
  \end{cases}
\end{equation}

Subgradient w.r.t. $b$:
\begin{equation}
  \partial_b J =
  \begin{cases}
    0 & \text{if } m_i \geq 1 \\
    -\dfrac{C}{n}\,y_i & \text{if } m_i < 1
  \end{cases}
\end{equation}

\subsection{Pegasos Learning-Rate Schedule}

The Pegasos algorithm~\cite{shalev2007} uses $\lambda = 1/C$ and schedule:
\begin{equation}
  \eta_t = \frac{1}{\lambda \cdot t}
\end{equation}
ensuring convergence in $O(1/(\lambda\varepsilon))$ iterations to $\varepsilon$-accuracy.

\subsection{RBF Kernel SVM}

The dual formulation with kernel $K$:
\begin{equation}
  \max_{\boldsymbol{\alpha}} \sum_i \alpha_i - \frac{1}{2}\sum_{i,j}
  \alpha_i\alpha_j y_i y_j K(\mathbf{x}_i,\mathbf{x}_j)
\end{equation}
\begin{equation}
  \text{s.t.}\quad 0 \leq \alpha_i \leq C,\quad \sum_i \alpha_i y_i = 0
\end{equation}

RBF kernel: $K(\mathbf{x}_i, \mathbf{x}_j) = \exp(-\gamma\|\mathbf{x}_i - \mathbf{x}_j\|^2)$.

\subsection{Probability Calibration (Platt Scaling)}

After training, the calibrated probability is:
\begin{equation}
  P(y{=}1 \mid f(\mathbf{x})) = \frac{1}{1 + \exp(A\cdot f(\mathbf{x}) + B)}
\end{equation}
where $A, B$ are fitted by minimising the cross-entropy on a held-out calibration set.

The Expected Calibration Error (ECE) is:
\begin{equation}
  \text{ECE} = \sum_{b=1}^B \frac{|B_b|}{n}
  \left|\overline{\text{acc}}(B_b) - \overline{\text{conf}}(B_b)\right|
\end{equation}

% ─────────────────────────────────────────────
\section{Methodology}
\label{sec:methodology}

\subsection{Dataset}
We use the ASVspoof~2019 Logical Access (LA) partition following the official
train/dev/eval split. Training: 2580 genuine + 22800 spoof. The \emph{eval}
partition is treated as the held-out internal test set. For external evaluation,
we use the In-the-Wild dataset~\cite{muller2022}.

\subsection{Audio Preprocessing}
All audio is decoded to 16~kHz mono, peak-normalised, and segmented into 3-second
windows with 1.5-second hop (50\% overlap). Segments shorter than 0.5~seconds are
rejected. Quality flags are raised for clipping $>1\%$ or speech ratio $<10\%$.

\subsection{Feature Extraction}
Feature vector construction:
\begin{itemize}
  \item \textbf{MFCC}: 20 coefficients + $\Delta$ + $\Delta\Delta$ (with 6 statistics each).
  \item \textbf{LFCC}: 20 coefficients + $\Delta$ + $\Delta\Delta$ using a linear filterbank.
  \item \textbf{Spectral}: centroid, bandwidth, rolloff, flatness, 7-band contrast, ZCR.
  \item \textbf{Prosodic}: RMS energy, energy variation, F0 (pyin), voiced ratio.
\end{itemize}
Statistics computed per frame sequence: mean, std, median, 25th/75th percentile, max.

\subsection{Models Compared}
\begin{enumerate}
  \item From-scratch Linear SVM (mini-batch subgradient)
  \item scikit-learn LinearSVC
  \item scikit-learn SVC(linear kernel)
  \item scikit-learn SVC(RBF kernel) [\textbf{deployed}]
  \item Logistic Regression (baseline)
  \item Random Forest (baseline)
\end{enumerate}

Hyperparameter grid: $C \in \{0.01, 0.1, 1, 10, 100\}$;
$\gamma \in \{\text{scale}, 0.001, 0.01, 0.1\}$ (RBF only).
3-fold stratified CV on training partition.

% ─────────────────────────────────────────────
\section{Results and Discussion}
\label{sec:results}

\subsection{Feature Ablation Study}

\begin{table}[htbp]
\caption{Feature Ablation — Internal Dev Set Performance}
\centering
\begin{tabular}{lcccc}
\toprule
Feature Set & Accuracy & F1 & AUC & EER \\
\midrule
MFCC only            & [pending] & [pending] & [pending] & [pending] \\
LFCC only            & [pending] & [pending] & [pending] & [pending] \\
MFCC + LFCC          & [pending] & [pending] & [pending] & [pending] \\
MFCC + LFCC + Spec.  & [pending] & [pending] & [pending] & [pending] \\
All features         & [pending] & [pending] & [pending] & [pending] \\
\bottomrule
\end{tabular}
\label{tab:ablation}
\end{table}

Results in Table~\ref{tab:ablation} are pending dataset download and
training execution. Run the Colab notebook to populate these values.

\subsection{Model Comparison}

\begin{table}[htbp]
\caption{Model Comparison — Internal Test Set}
\centering
\begin{tabular}{lcccc}
\toprule
Model & Acc. (\%) & F1 & AUC & EER (\%) \\
\midrule
Scratch Linear SVM   & [pending] & --- & --- & --- \\
sklearn LinearSVC    & [pending] & --- & --- & --- \\
SVC (linear)         & [pending] & --- & --- & --- \\
SVC (RBF)            & [pending] & --- & --- & --- \\
Logistic Regression  & [pending] & --- & --- & --- \\
Random Forest        & [pending] & --- & --- & --- \\
\bottomrule
\end{tabular}
\label{tab:model_comparison}
\end{table}

\subsection{Generalization Analysis}

Internal vs. external dataset performance gap will be reported here after
running the full experiment (Cells 29-30 in the Colab notebook).

\subsection{Calibration Quality}

Pre- and post-Platt-scaling Brier scores and ECE values are logged in
Cell~27 of the notebook and saved to \texttt{results/metrics/}.

% ─────────────────────────────────────────────
\section{Limitations}
\label{sec:limitations}

\begin{enumerate}
  \item \textbf{Generator generalisation}: Novel TTS/VC systems not seen during
    training may evade the detector.
  \item \textbf{Domain shift}: Performance degrades with different recording
    conditions, acoustic environments, or heavy codec compression.
  \item \textbf{Short audio}: Recordings $<0.5$~s are rejected; $<1$~s predictions
    are less reliable.
  \item \textbf{Colab hosting}: The Gradio app depends on an active Colab session.
  \item \textbf{Not forensic evidence}: This system should not be sole evidence
    in consequential decisions.
\end{enumerate}

% ─────────────────────────────────────────────
\section{Conclusion}
\label{sec:conclusion}

We have presented a complete, academically rigorous deepfake audio detection
pipeline built on Support Vector Machines implemented from first principles.
The mathematical treatment covers primal and dual SVM formulations, hinge loss,
subgradient computation, Pegasos learning-rate schedules, the kernel trick, and
Platt probability calibration. The interactive EchoVerify AI web application
enables real-time prediction on user-uploaded audio. Future work should explore
end-to-end neural approaches (e.g., Wav2Vec-based CM systems), multi-condition
training for improved domain robustness, and ensemble calibration methods.

% ─────────────────────────────────────────────
\section{Individual Contribution Statement}
\begin{itemize}
  \item \textbf{[Author 1]}: [Describe specific contributions, e.g., SVM
    implementation, feature extraction, evaluation]
  \item \textbf{[Author 2]}: [Describe specific contributions, e.g., dataset
    processing, Gradio app, report writing]
\end{itemize}

% ─────────────────────────────────────────────
\begin{thebibliography}{99}

\bibitem{wavenet2016}
A. van den Oord et al., ``WaveNet: A Generative Model for Raw Audio,''
\emph{arXiv:1609.03499}, 2016.

\bibitem{tacotron2017}
Y. Wang et al., ``Tacotron: Towards End-to-End Speech Synthesis,''
in \emph{Proc. Interspeech}, 2017, pp. 4006--4010.

\bibitem{valle2023}
C. Wang et al., ``Neural Codec Language Models are Zero-Shot Text to Speech
Synthesizers,'' \emph{arXiv:2301.02111}, 2023.

\bibitem{asvspoof2019}
X. Wang et al., ``ASVspoof 2019: A Large-Scale Public Database of Synthesized,
Converted and Replayed Speech,'' \emph{Computer Speech \& Language},
vol. 64, 2020.

\bibitem{sahidullah2015}
Md. Sahidullah et al., ``A Comparison of Features for Synthetic Speech Detection,''
in \emph{Proc. Interspeech}, 2015, pp. 2087--2091.

\bibitem{todisco2016}
M. Todisco et al., ``A New Feature for Automatic Speaker Verification Anti-Spoofing:
Constant Q Cepstral Coefficients,'' in \emph{Proc. Odyssey}, 2016, pp. 283--290.

\bibitem{platt1999}
J. Platt, ``Probabilistic Outputs for Support Vector Machines and Comparisons
to Regularized Likelihood Methods,'' in \emph{Advances in Large Margin Classifiers},
MIT Press, 1999, pp. 61--74.

\bibitem{niculescu2005}
A. Niculescu-Mizil and R. Caruana, ``Predicting Good Probabilities with Supervised
Learning,'' in \emph{Proc. ICML}, 2005, pp. 625--632.

\bibitem{shalev2007}
S. Shalev-Shwartz et al., ``Pegasos: Primal Estimated sub-GrAdient SOlver for SVM,''
in \emph{Proc. ICML}, 2007, pp. 807--814.

\bibitem{jung2022aasist}
J. Jung et al., ``AASIST: Audio Anti-Spoofing using Integrated Spectro-Temporal Graph
Attention Networks,'' in \emph{Proc. ICASSP}, 2022.

\bibitem{muller2022}
N. M\"{u}ller et al., ``Does Audio Deepfake Detection Generalize?,''
in \emph{Proc. Interspeech}, 2022.

\end{thebibliography}

\end{document}
